\documentclass[12pt]{article}
\usepackage{amsmath,amssymb}
\usepackage[margin=1in]{geometry}
\pagestyle{empty}
\begin{document}

\begin{center}
{\Large\bfseries The Price of Memory}
\end{center}

\medskip

Two numbers govern the growth of sequences built from binary choices.

Without memory---each choice independent of the last---the growth rate is the \textbf{golden ratio}:
\[
\varphi = \frac{1+\sqrt{5}}{2} = 1.61803\ldots \qquad (\text{Fibonacci: } \lambda^2 = \lambda + 1)
\]

With one step of memory---each choice sees what happened just before---the growth rate is the \textbf{tribonacci constant}:
\[
\tau = 1.83929\ldots \qquad (\text{tribonacci: } \lambda^3 = \lambda^2 + \lambda + 1)
\]

The difference
\[
\tau - \varphi = 0.22125\ldots
\]
is the \textbf{price of memory}: the exact cost, in growth rate per step, of remembering what just happened.

Tournaments pay this price. Every comparison ``A beats B'' echoes into the next: not as a new decision, but as a \emph{record}---the third state (besides win and lose) that makes the $3 \times 3$ transfer matrix irreducible, the recurrence cubic instead of quadratic, the golden ratio not quite enough.

The entire theory of pairwise comparison lives in that gap of $0.221$.

\end{document}
